\documentclass[conference]{IEEEtran}
\IEEEoverridecommandlockouts
\usepackage[utf8]{inputenc}
\usepackage[T1]{fontenc}
\usepackage{amsmath,amssymb,amsfonts}
\usepackage{booktabs}
\usepackage{url}
\usepackage{hyperref}
\hypersetup{colorlinks=true,linkcolor=black,citecolor=black,urlcolor=blue}

\renewcommand{\abstractname}{Resumen}
\renewcommand{\IEEEkeywordsname}{Palabras clave}
\renewcommand{\tablename}{TABLA}
\renewcommand{\refname}{Referencias}

\newcommand{\Fdos}{\mathbb{F}_2}
\newcommand{\xor}{\oplus}

\begin{document}

\title{Criptoanálisis lineal y diferencial de una variante ligera de Keccak:\\
criterios de diseño verificados y refutación de la premisa de difusión}

\author{\IEEEauthorblockN{Oscar Toledo}
\IEEEauthorblockA{Maestría en Ciencias de la Computación\\
Universidad Nacional de Ingeniería (UNI)\\
Lima, Perú}}

\maketitle

\begin{abstract}
Se aborda el rediseño de Keccak para hardware limitado siguiendo los cinco
requerimientos planteados, con una metodología basada en medición y verificación
formal antes que en afirmaciones de diseño. Los resultados contradicen la premisa
de partida: se demuestra que Keccak-f[200] alcanza el 80\,\% de difusión en dos
rondas y Keccak-f[800] en tres, de modo que el número elevado de rondas del
estándar no responde a necesidades de difusión sino a margen de seguridad. Se
evalúan seis capas lineales alternativas construidas solo con XOR y rotaciones;
ninguna supera la cota diferencial de $\theta$, y se identifica la razón
estructural: $\theta$ obtiene dependencia de once términos por bit con dos XOR
por bit amortizados, gracias a que reutiliza cinco veces cada paridad de columna.
Del proceso se derivan cuatro criterios de diseño verificados algebraicamente, y
se documenta una vulnerabilidad catastrófica de la capa no lineal parcial
propuesta para abaratar el circuito: admite un subespacio invariante de $2^{8}$
diferencias que atraviesan cualquier número de rondas con probabilidad uno. Todas
las cotas diferenciales se certifican con CP-SAT.
\end{abstract}

\begin{IEEEkeywords}
Keccak, SHA-3, criptoanálisis lineal, criptoanálisis diferencial, LAT, DDT,
CP-SAT, hardware ligero.
\end{IEEEkeywords}

\section{Criptoanálisis lineal preliminar}

\subsection{Aproximación afín de una función no lineal}

El criptoanálisis lineal aproxima una función no lineal mediante funciones afines
y explota el sesgo resultante. Para $x \lor y$ sobre dos variables, las cuatro
aproximaciones consideradas coinciden con la función en tres de los cuatro puntos
(Tabla~\ref{tab:or}), cada una fallando en uno distinto: la constante $1$ falla en
$(0,0)$, $x$ en $(0,1)$, $y$ en $(1,0)$ y $x \xor y$ en $(1,1)$.

\begin{table}[!ht]
\caption{Aproximaciones afines de $x \lor y$}
\label{tab:or}
\centering
\begin{tabular}{@{}ccccccc@{}}
\toprule
$x$ & $y$ & $x \lor y$ & $1$ & $x$ & $y$ & $x \xor y$ \\
\midrule
0 & 0 & 0 & 1 & 0 & 0 & 0 \\
0 & 1 & 1 & 1 & 0 & 1 & 1 \\
1 & 0 & 1 & 1 & 1 & 0 & 1 \\
1 & 1 & 1 & 1 & 1 & 1 & 0 \\
\midrule
\multicolumn{3}{c}{Probabilidad} & $3/4$ & $3/4$ & $3/4$ & $3/4$ \\
\bottomrule
\end{tabular}
\end{table}

El valor $3/4$ es óptimo. Sobre dos variables existen ocho funciones afines; las
cuatro restantes ($0$, $x\xor1$, $y\xor1$, $x\xor y\xor1$) aciertan $1/4$ o $2/4$.
Se dice entonces que $x \lor y$ tiene \emph{no linealidad} $1$: dista una única
posición de la función afín más próxima. El sesgo es
$\varepsilon = |3/4 - 1/2| = 1/4$, y es la magnitud que hace explotable la
aproximación.

\subsection{Aproximación sobre la caja-S de cuatro bits}

Para la caja-S dada por $S(x) = \mathtt{c2\,6f\,90\,4d\,3e\,b8\,a7\,15}$ se evalúa
la aproximación $y_3 = x_2 \xor x_3$, con $x_1$ como bit más significativo. El
recuento exhaustivo sobre los dieciséis valores de entrada arroja
\begin{equation}
\Pr[y_3 = x_2 \xor x_3] = \frac{10}{16} = 0{,}625,
\qquad \varepsilon = \frac{1}{8}.
\end{equation}
Por la regla de Matsui, un ataque lineal que explote esta aproximación requiere del
orden de $\varepsilon^{-2} = 64$ textos claros conocidos.

\subsection{Tabla de aproximaciones lineales}

La LAT resume todas las aproximaciones posibles. Cada casilla es
\begin{equation}
\mathrm{LAT}[\alpha][\beta] =
\#\{x : \alpha \cdot x = \beta \cdot S(x)\} - 8,
\end{equation}
donde $\alpha \cdot x$ denota el XOR de los bits de $x$ seleccionados por la
máscara $\alpha$. De ahí se obtiene $p = (8 + \mathrm{LAT})/16$ y el sesgo
$|\mathrm{LAT}|/16$. La Tabla~\ref{tab:lat} recoge seis casillas representativas.

\begin{table}[!ht]
\caption{Seis valores representativos de la LAT}
\label{tab:lat}
\centering
\begin{tabular}{@{}cccccl@{}}
\toprule
$\alpha$ & $\beta$ & LAT & $p$ & sesgo & Lectura \\
\midrule
0 & 0 & $+8$ & $1$      & $1/2$ & referencia \\
0 & 1 & $0$  & $1/2$    & $0$   & fila nula: $S$ balanceada \\
1 & 0 & $0$  & $1/2$    & $0$   & columna nula: $S$ biyectiva \\
6 & 2 & $+2$ & $10/16$  & $1/8$ & la aproximación anterior \\
7 & 3 & $+6$ & $14/16$  & $3/8$ & sesgo máximo positivo \\
1 & 5 & $-6$ & $2/16$   & $3/8$ & sesgo máximo negativo \\
\bottomrule
\end{tabular}
\end{table}

La casilla $(6,2)$ vale $+2$, que es exactamente $10-8$: la aproximación estudiada
en el apartado anterior es una de las $256$ casillas de la tabla. Conviene señalar
que esta caja-S es débil: su $|\mathrm{LAT}|$ máximo es $6$ y su DDT máximo es $8$,
mientras que una caja-S de cuatro bits óptima alcanza $4$ en ambas métricas. Un
DDT de $8$ sobre $16$ implica la existencia de una transición diferencial de
probabilidad $1/2$.

\section{Modificación estructural e inyección dinámica}

\subsection{Estructura propuesta}

La variante sustituye el estado tridimensional $5\times5\times w$ por una
disposición en rodajas de cinco carriles de $w$ bits, con $5w$ bits en total. Cada
ronda aplica una capa lineal $\Lambda$ dentro de cada carril y la capa no lineal
$\chi$ sobre columnas de cinco bits, una por índice $z$. La difusión entre
carriles proviene exclusivamente de $\chi$ y la difusión a lo largo del carril
exclusivamente de las rotaciones.

\subsection{Las dos variables dinámicas}

Se inyectan dos parámetros en tiempo de ejecución:
\begin{itemize}
\item $s$, \textbf{nivel de seguridad}, que determina el número de rondas $R(s)$ y
la capacidad $c(s)$ de la esponja;
\item $d$, \textbf{dominio de aplicación}, que se inyecta como constante de
separación de dominio en el relleno.
\end{itemize}

La interacción se diseña para no introducir debilidades. La separación de dominio
mediante constantes es el mecanismo estándar de SHA-3, que distingue así SHA-3 de
SHAKE mediante sufijos distintos; al ser una constante pública que solo desplaza
el punto de partida, no altera las propiedades diferenciales ni lineales.

El punto delicado es $s$. Si el número de rondas dependiera de un valor
correlacionado con un secreto, el tiempo de cómputo revelaría información sobre
ese secreto: un canal lateral temporal directo. El requisito de diseño es por
tanto que $s$ sea un parámetro \emph{público}, fijado por la aplicación al inicio
de la sesión y constante durante ella. Esta condición no es una formalidad: en un
trabajo previo se analizó una variante cuyo número de rondas dependía de un
contador de intentos fallidos, y esa construcción sí filtra el estado de
autenticación a través del tiempo de ejecución.

\section{Rediseño de la capa lineal}

\subsection{Refutación de la premisa}

El planteamiento del caso afirma que Keccak-p requiere demasiadas rondas para
asegurar difusión completa. La afirmación es contrastable y resulta falsa.

Se mide la difusión mediante la matriz de dependencia sobre $\Fdos$: se define
$D_r[i][j] = 1$ si el bit de salida $i$ depende del bit de entrada $j$ tras $r$
rondas, se compone ronda a ronda mediante el producto booleano, y se toma como
difusión la densidad de la matriz. Cada paso aporta sus dependencias exactas:
$\theta$ once entradas por bit, $\rho$ y $\pi$ una (son permutaciones) y $\chi$
tres.

\begin{table}[t]
\caption{Rondas hasta alcanzar el 80\,\% de difusión}
\label{tab:dif}
\centering
\begin{tabular}{@{}lcc@{}}
\toprule
Construcción & Estado & Rondas \\
\midrule
Keccak-f[200]         & 200 bits & \textbf{2} \\
Keccak-f[400]         & 400 bits & \textbf{2} \\
Keccak-f[800]         & 800 bits & \textbf{3} \\
Variante, $\Sigma$ por carril    & 200 bits & 3 \\
Variante, paridad $+$ $\Sigma$   & 200 bits & 3 \\
\bottomrule
\end{tabular}
\end{table}

La Tabla~\ref{tab:dif} muestra que Keccak alcanza difusión prácticamente completa
en dos o tres rondas. Las veinticuatro rondas del estándar no responden a
necesidades de difusión, sino a margen de seguridad frente a criptoanálisis
diferencial y algebraico. En consecuencia, el objetivo de lograr el 80\,\% en un
número significativamente menor de rondas carece de recorrido: solo queda una
ronda, y para que un bit dependiera del 80\,\% de los $200$ del estado en un único
paso harían falta unos $160$ XOR por bit, frente a los dos por bit de $\theta$.

\subsection{Por qué $\theta$ es difícil de superar}

$\theta$ calcula cinco paridades de columna y las \emph{reutiliza en los
veinticinco carriles}. Ese factor de reutilización de cinco explica que consiga
dependencia de once términos por bit con solo dos XOR por bit amortizados. Una
estructura de cinco carriles no puede replicarlo, porque hay menos consumidores
por cada paridad calculada. Se trata de una limitación estructural y no de una
insuficiencia de la búsqueda realizada.

\subsection{Criterios de invertibilidad}

Durante la exploración de alternativas se descartaron varios diseños por no ser
permutaciones, requisito indispensable en una construcción de esponja. De los
fallos se extraen dos criterios verificados algebraicamente.

\textbf{Criterio 1: el número de términos por carril debe ser impar.} El mapa
$L \mapsto L \xor \mathrm{rot}(L,a)$ es siempre singular, porque el vector
todo-unos permanece invariante bajo rotación y se cancela consigo mismo. En
general, con $k$ términos el estado todo-unos se transforma en el XOR de $k$
copias de sí mismo, que es nulo cuando $k$ es par. La verificación numérica lo
confirma: con cuatro términos por carril el modelo devuelve cero cajas-S activas,
señal inequívoca de singularidad.

\textbf{Criterio 2: la paridad debe cubrir cuatro carriles, no cinco.} Al añadir
la paridad de los cinco carriles a los tres términos propios se obtienen cuatro
términos, y reaparece el problema anterior con deficiencia de rango exactamente
uno. Restringir la paridad a cuatro carriles restablece la invertibilidad.

\section{Rediseño de la capa no lineal}

\subsection{La propuesta y su motivación}

El paso $\chi$ cuesta, por bit, una compuerta ANDN y un XOR, con profundidad dos.
Reducir ese coste directamente es difícil, pero sí puede reducirse el
\emph{número} de compuertas AND aplicando $\chi$ solo a un subconjunto de las
columnas en cada ronda, alternando el subconjunto entre rondas. La motivación es
sólida: en implementaciones enmascaradas contra canal lateral el coste está
dominado por las compuertas AND, ya que los XOR son lineales y las máscaras los
atraviesan sin refresco. Reducir a la mitad las AND reduce a la mitad el coste del
enmascaramiento.

\subsection{Demostración de que la propuesta falla}

La propuesta es insegura, y la demostración es constructiva. Sea $\Lambda$ la capa
lineal y $P_r$ la proyección sobre las columnas que llevan $\chi$ en la ronda $r$.
Si una diferencia $d$ cumple
\begin{equation}\label{eq:nucleo}
P_r\bigl(\Lambda^{r+1} d\bigr) = 0, \qquad r = 0,\dots,R-1,
\end{equation}
entonces $\chi$ recibe cero en todas las columnas activas y devuelve cero, de modo
que la ronda completa es lineal para esa diferencia y la trayectoria tiene
probabilidad uno. El conjunto de tales $d$ es un subespacio vectorial, el núcleo
del sistema~\eqref{eq:nucleo}, y su dimensión se calcula por eliminación gaussiana
sobre $\Fdos$.

\begin{table}[t]
\caption{Dimensión del subespacio de diferencias que esquivan toda no linealidad}
\label{tab:nucleo}
\centering
\begin{tabular}{@{}lccccc@{}}
\toprule
Capa no lineal & $R{=}1$ & $R{=}2$ & $R{=}3$ & $R{=}4$ & $R{=}6$ \\
\midrule
$\chi$ completo           & $0$      & $0$    & $0$   & $0$   & $0$ \\
$\chi$ parcial (1/2)      & $100$    & $8$    & $8$   & $8$   & $8$ \\
$\chi$ parcial (1/4)      & $150$    & $101$  & $59$  & --    & -- \\
\bottomrule
\end{tabular}
\end{table}

El resultado es concluyente. La sucesión $V_1 \supseteq V_2 \supseteq \cdots$ es
decreciente y \emph{se estabiliza} en el mayor subespacio invariante bajo
$\Lambda$ contenido en la región inactiva. Para $\chi$ parcial a la mitad, esa
estabilización ocurre en dimensión $8$: existen $2^{8} = 256$ diferencias no nulas
que atraviesan \emph{cualquier} número de rondas con probabilidad uno. Añadir
rondas no repara el defecto.

El hallazgo se confirmó por tres vías independientes: el modelo CP-SAT devuelve
cero cajas-S activas, la propagación directa de la diferencia sobre el estado lo
reproduce, y el cálculo del núcleo lo explica y cuantifica.

\subsection{Ecuaciones correctas}

\textbf{Criterio 3: el núcleo debe ser trivial.} La condición de diseño es que el
sistema~\eqref{eq:nucleo} solo admita la solución nula. Con $\chi$ aplicado a todas
las columnas el núcleo es trivial para todo $R$ (Tabla~\ref{tab:nucleo}), y el
modelo CP-SAT confirma cotas diferenciales no nulas. La ecuación correcta de la
capa no lineal es, por tanto, la aplicación completa
\begin{equation}
\chi_i = x_i \xor \bigl(\overline{x_{i+1}} \cdot x_{i+2}\bigr),
\qquad i = 0,\dots,4,
\end{equation}
sobre todas las columnas, y el ahorro de compuertas AND debe buscarse por otra vía.
Se ensayó también añadir mezcla entre carriles a la capa lineal para reparar la
versión parcial; el núcleo se reduce de $2^{8}$ a $2^{4}$ pero no desaparece, de
modo que la corrección es insuficiente.

\section{Demostración formal mediante MILP}

\subsection{Modelo}

Se modela la propagación de diferencias asignando una variable binaria a cada bit
del estado. La capa lineal se impone mediante restricciones XOR exactas; la capa
no lineal, exigiendo que el par (diferencia de entrada, diferencia de salida) de
cada caja-S sea una transición válida de la DDT de $\chi$, que consta de $317$
transiciones. El objetivo es minimizar el número de cajas-S activas.

La resolución se realiza con CP-SAT, que trata de forma nativa el XOR y las tablas
de transiciones. En trabajo previo se comprobó que la formulación MILP equivalente
no cierra el gap para $R \ge 2$, porque cerca del 90\,\% de sus variables
corresponden a la codificación \emph{big-M} de la DDT y las restricciones de
paridad sobre $\Fdos$ se relajan muy débilmente.

\subsection{Comparación con la línea base}

\begin{table}[t]
\caption{Mínimo de cajas-S activas (200 bits, 40 cajas por ronda)}
\label{tab:milp}
\centering
\begin{tabular}{@{}lccc@{}}
\toprule
Construcción & $R{=}1$ & $R{=}2$ & $R{=}3$ \\
\midrule
Keccak-f[200] (línea base) & $1$ & $4$ & $\mathbf{10}$ \\
Variante, $k=3$ términos   & $1$ & $4$ & $7$ \\
\bottomrule
\end{tabular}
\end{table}

Todos los valores de la Tabla~\ref{tab:milp} están certificados: CP-SAT devuelve
\texttt{OPTIMAL} con gap cerrado. La variante iguala a la línea base hasta dos
rondas y queda por debajo en tres. La razón es la difusión local: con $k$ términos
por carril, un bit activo alcanza a lo sumo $k$ columnas, de modo que la
trayectoria mínima observada es $1 \to 3 \to 3$, mientras $\theta$, con once
términos por bit, fuerza $1 \to 3 \to 6$.

Para la variante con $k=3$ se obtiene además una cota inferior demostrada de $10$
cajas activas en cuatro rondas. El intercambio queda así cuantificado con
precisión: \emph{la variante alcanza la garantía que Keccak ofrece en tres rondas
empleando cuatro}.

\textbf{Criterio 4: la cota diferencial la fija el número de términos, no las
constantes.} Se realizó una búsqueda sobre $17\,125$ conjuntos de constantes de
rotación, filtrados por dispersión máxima de los desplazamientos relativos,
invertibilidad y velocidad de difusión. Los mejores candidatos según ese criterio
arrojaron cotas de $5$ cajas activas en tres rondas, \emph{peores} que las $7$ de
partida, y el valor en dos rondas resultó ser $4$ para todos los conjuntos
probados. La conclusión, contraria a la intuición inicial, es que la elección fina
de constantes no gobierna la cota diferencial en esta familia.

\section{Justificación de la variante}

\subsection{Balance}

\begin{table}[t]
\caption{Comparación estructural a igualdad de estado}
\label{tab:balance}
\centering
\begin{tabular}{@{}lcc@{}}
\toprule
 & Keccak-f[200] & Variante \\
\midrule
Carriles                        & 25   & \textbf{5} \\
Rotaciones distintas            & 25   & \textbf{10} \\
Permutación $\pi$               & sí   & \textbf{no} \\
XOR por ronda                   & 400  & 400 \\
Rondas para $\ge 10$ cajas      & \textbf{3} & 4 \\
\bottomrule
\end{tabular}
\end{table}

La variante no mejora a Keccak en compuertas: para igualar la garantía necesita
una ronda adicional, lo que supone $1600$ XOR frente a $1200$. Su ventaja es
estructural: cinco carriles en lugar de veinticinco, diez rotaciones en lugar de
veinticinco y ausencia de la permutación $\pi$. En implementaciones compactas
serializadas, donde el enrutado y el número de multiplexores dominan el área, esa
simplicidad puede compensar la ronda extra; afirmarlo con rotundidad exigiría una
síntesis sobre tecnología concreta, que queda fuera de este trabajo.

\subsection{Las dos vulnerabilidades de la propuesta}

\textbf{Primera: la capa no lineal parcial admite diferenciales de probabilidad
uno.} Es la vulnerabilidad grave, demostrada en la Sección IV. El subespacio
invariante de dimensión $8$ permite atravesar cualquier número de rondas sin
activar una sola caja-S, lo que anula por completo la resistencia diferencial. La
corrección consiste en aplicar $\chi$ a todas las columnas, con la verificación de
que el núcleo del sistema~\eqref{eq:nucleo} es trivial.

\textbf{Segunda: la difusión local degrada la cota diferencial.} Al confinar la
capa lineal dentro de cada carril, la variante activa siete cajas-S en tres rondas
frente a las diez de la línea base. Un atacante dispone entonces de trayectorias
de probabilidad $2^{-14}$ donde en Keccak tendría $2^{-20}$. La corrección
consiste en compensar con una ronda adicional, con el coste correspondiente, o en
introducir mezcla entre carriles en la capa lineal, respetando los criterios 1
y 2 de invertibilidad.

\section{Conclusiones}

\begin{enumerate}
\item La premisa del caso no se sostiene. Keccak alcanza el 80\,\% de difusión en
dos o tres rondas según el tamaño de estado; sus veinticuatro rondas responden a
margen de seguridad, no a difusión.
\item Ninguna de las seis capas lineales evaluadas supera la cota diferencial de
$\theta$. La razón es estructural: $\theta$ reutiliza cinco veces cada paridad de
columna, y una disposición de cinco carriles no puede igualar ese factor.
\item La capa no lineal parcial, propuesta para abaratar el circuito, es
insegura: admite $2^{8}$ diferenciales de probabilidad uno que ningún número de
rondas elimina.
\item Se derivan cuatro criterios de diseño verificados: número impar de términos
por carril, paridad sobre cuatro carriles, núcleo trivial, y dependencia de la
cota diferencial respecto del número de términos y no de las constantes.
\item La variante iguala la garantía diferencial de Keccak empleando una ronda
adicional, a cambio de una estructura sensiblemente más simple.
\end{enumerate}

Como trabajo futuro, quedan pendientes la certificación de la variante con cinco
términos por carril en tres rondas ---donde se halló una trayectoria de catorce
cajas activas sin poder demostrar su minimalidad--- y la síntesis sobre tecnología
concreta que permitiría cuantificar la ventaja de área.

\begin{thebibliography}{9}
\bibitem{fips202} National Institute of Standards and Technology, \emph{SHA-3
Standard: Permutation-Based Hash and Extendable-Output Functions}, FIPS PUB 202,
2015.
\bibitem{keccak} G. Bertoni, J. Daemen, M. Peeters y G. Van Assche, \emph{The
Keccak Reference}, versión 3.0, 2011.
\bibitem{matsui} M. Matsui, ``Linear cryptanalysis method for DES cipher,'' en
\emph{EUROCRYPT}, 1993, pp. 386--397.
\bibitem{biham} E. Biham y A. Shamir, ``Differential cryptanalysis of DES-like
cryptosystems,'' \emph{Journal of Cryptology}, vol. 4, pp. 3--72, 1991.
\bibitem{ascon} C. Dobraunig, M. Eichlseder, F. Mendel y M. Schläffer,
\emph{Ascon v1.2}, sumisión al concurso NIST Lightweight Cryptography, 2021.
\bibitem{ortools} L. Perron y F. Didier, \emph{CP-SAT}, Google OR-Tools.
\end{thebibliography}

\end{document}
